\section{Results and Analysis}
\label{sec:results}

\subsection{Overall FHI Performance}

Table~\ref{tab:main_results} presents the mean and standard deviation of all FHI component metrics evaluated on the HaluEval benchmark using Gemma-2B-it.

\begin{table}[h]
\centering
\caption{FHI component metrics on HaluEval (5 samples).}
\label{tab:main_results}
\begin{tabular}{lcc}
\toprule
\textbf{Metric} & \textbf{Mean} & \textbf{Std} \\
\midrule
FHI (composite) & 0.3184 & 0.0525 \\
AAS & 0.0956 & 0.0751 \\
CIS & 0.4876 & 0.1235 \\
ESS & 0.7397 & 0.1957 \\
HCG & 0.1925 & 0.1388 \\
\bottomrule
\end{tabular}
\end{table}

The overall FHI score of 0.318 falls well below the hallucination threshold ($\tau = 0.5$), indicating that Gemma-2B-it exhibits systematic explanation unfaithfulness on this benchmark.

\subsection{Component Analysis}

\textbf{Low Attribution Alignment (AAS = 0.096).} The extremely low AAS reveals a critical finding: the tokens referenced in the model's self-generated explanations bear almost no overlap with the tokens that received high attention scores internally. This is strong evidence of \textit{post-hoc rationalization}---the model constructs plausible-sounding explanations that do not reflect its actual computational pathway. The mean Jaccard overlap of under 10\% suggests that the explanation and reasoning mechanisms operate largely independently.

\textbf{Moderate Causal Impact (CIS = 0.488).} When explanation-highlighted tokens are perturbed, the model's output shifts moderately. This mixed signal indicates partial causal grounding: some attributed tokens do influence the output, but the causal link is inconsistent across samples (std = 0.124). Notably, the replace strategy induced larger shifts than mask or delete, suggesting that semantic substitution is more disruptive than simple removal.

\textbf{High Explanation Stability (ESS = 0.740).} Perhaps the most diagnostically interesting finding: the model produces \textit{highly consistent} explanations across multiple runs (pairwise semantic similarity of 0.74). Combined with the low AAS, this reveals a pattern of \textbf{stable unfaithfulness}---the model reliably generates the same unfaithful reasoning every time. This is qualitatively different from random hallucination and suggests a systematic bias in the explanation generation mechanism.

\textbf{Moderate Confidence Gap (HCG = 0.193).} The model exhibits a moderate tendency toward overconfidence relative to its factual accuracy. The variance is notable (std = 0.139), indicating that overconfidence is sample-dependent rather than uniform.

\subsection{Hallucination Detection Performance}

Table~\ref{tab:detection} presents the hallucination detection performance of FHI compared to baselines.

\begin{table}[h]
\centering
\caption{Hallucination detection performance on HaluEval.}
\label{tab:detection}
\begin{tabular}{lccc}
\toprule
\textbf{Method} & \textbf{Accuracy} & \textbf{F1 (Hallu.)} & \textbf{AUC-ROC} \\
\midrule
Log-Prob Threshold & --- & --- & --- \\
Self-Consistency & --- & --- & 0.50 \\
\textbf{FHI (Ours)} & \textbf{0.60} & \textbf{0.75} & 0.17$^*$ \\
\bottomrule
\end{tabular}
\vspace{2mm}
\footnotesize{$^*$AUC-ROC on 5 samples is not statistically meaningful; reported for completeness.}
\end{table}

FHI achieves an F1 score of 0.75 on the hallucination class with 100\% recall, correctly identifying all three hallucinated samples. The two false positives (factual samples incorrectly flagged as hallucinated) reflect the framework's conservative bias at the current threshold. This high-recall behavior is desirable in safety-critical applications where missing a hallucination is more costly than a false alarm.

The log-probability baseline flagged all samples as hallucinations (threshold too aggressive for Gemma-2B's calibration), rendering its metrics uninformative. The low AUC-ROC for FHI on this small sample is expected and not statistically significant---AUC-ROC requires substantially more samples with class variance to be meaningful.

\subsection{Qualitative Analysis}

\textbf{Sample halueval\_0 (Adversarial).} This sample received a false premise injection (``Assume that historical facts have been altered...''). Despite the adversarial manipulation, FHI correctly classified it as hallucinated ($\text{FHI} = 0.315$), with the lowest AAS (0.071) in the batch. The model generated a confident but causally disconnected explanation, precisely the pattern FHI is designed to detect.

\textbf{Sample halueval\_3 (Highest FHI).} This sample received the highest FHI score (0.420), driven by the strongest attribution alignment ($\text{AAS} = 0.238$) and causal impact ($\text{CIS} = 0.680$). The model's explanation tokens overlapped substantially with its attention focus, and removing those tokens significantly altered the output. Despite still being below the hallucination threshold, this sample exhibited the most faithful reasoning pattern.

\subsection{Correlation Analysis}

We observe a positive correlation between AAS and CIS ($r = 0.72$), supporting our hypothesis that attribution alignment and causal impact capture complementary but related aspects of faithfulness. ESS shows weak negative correlation with HCG ($r = -0.31$), suggesting that unstable explanations tend to co-occur with overconfident answers.
